

%%% fhc mu-like

\begin{table}[htb]
\centering
\begin{tabular}{| c | c | c |}
\hline
  {\bf Source of uncertainty}    & \multicolumn{1}{c|}{{\bf ${\delta}N_{SK}$/$N_{SK}$}}\\
                                 & \multicolumn{1}{c|}{(Asimov B)}\\

\hline
\hline  
SKDet+FSI+SI                   & 4.03\%                        \\
SKDet only                   & 3.77\%                        \\
FSI+SI only                   & 1.43\%                        \\
\hline
  Flux                           & 3.58\%                        \\
\hline
\hline
  2p-2h (corr)                     & 3.53\%                        \\
  2p-2h-bar (corr)                  & 0.20\%                        \\
  NC other (uncorr)              & 0.76\%                        \\
  NC 1gamma (uncorr)             & 0.00\%                        \\
  XSec nue/numu (uncorr)                 & 0.01\%                        \\
\hline  
  XSec Tot (corr)                  & 4.01\%                        \\
  XSec Tot                       & 4.08\%                        \\
\hline
\hline
  Flux+XSec (ND280 constrained)                      & 2.77\%                        \\
  Flux+XSec (All)                      				& 2.87\%                        \\
\hline
  Flux+XSec+SKDet+FSI+SI         & 4.94\%                        \\
\hline
\hline
  Oscillations                   & 0.09\%                        \\
\hline
\hline
  All                            & 4.98\%                        \\
\hline
\hline

\end{tabular}
\caption{
 Effect of 1$\sigma$ variations of the groups of systematics parameters on total number of
 1-ring muon-like events for neutrino mode obtained by performing 10k toys MC.
The systematic parameter values measured with the ND280 data 
shown in 
\cref{tab:banff_uncertainties_flux_fhc},
~\ref{tab:banff_uncertainties_flux_rhc}
~\ref{tab:banff_uncertainties_xsec},
the nominal Super-K systematic parameters shown in \cref{tab:banff_uncertainties_skdetfsi}
and the oscillation parameters of the Asimov data set B shown in \cref{tab:osc_params_joint_asimov}
are used.
Since the parameters $\delta_{CP}$, $\theta_{23}$ and $\Delta m^2_{32}$ ($\Delta m^2_{13}$) and mass hierarchy
will be measured,
only the uncertainty on the parameters 
$\sin^2 \theta_{13}$,
$\sin^2 \theta_{12}$
and 
$\Delta m^2_{21}$ 
are taken into account.
}                                                                                                                                                                                          
\label{tab:nsk_systematic_table_summary_joint_asimov_B_fhc_mulike_Run17b}
\end{table}




%%%%% fhc e-like


\begin{table}[htb]
\centering
\begin{tabular}{| c | c | c |}
\hline
  {\bf Source of uncertainty}    & \multicolumn{1}{c|}{{\bf ${\delta}N_{SK}$/$N_{SK}$}}\\
                                 & \multicolumn{1}{c|}{(Asimov B)}\\

\hline
\hline  
SKDet+FSI+SI                   & 3.74\%                        \\
SKDet only                   & 2.65\%                        \\
FSI+SI only                   & 2.64\%                        \\
\hline
  Flux                           & 3.64\%                        \\
\hline
\hline
  2p-2h (corr)                     & 3.88\%                        \\
  2p-2h-bar (corr)                  & 0.07\%                        \\
  NC other (uncorr)              & 0.21\%                        \\
  NC 1gamma (uncorr)             & 1.95\%                        \\
  XSec nue/numu (uncorr)                 & 2.52\%                        \\
\hline
  XSec Tot (corr)                  & 4.15\%                        \\
  XSec Tot                       & 5.24\%                        \\
\hline
\hline
  Flux+XSec (ND280 constrained)                      & 2.88\%                        \\
  Flux+XSec (All)                      & 4.30\%                        \\
\hline
\hline
  Flux+XSec+SKDet+FSI+SI         & 5.78\%                        \\
\hline
\hline
  Oscillations                   & 4.23\%                        \\
\hline
\hline
  All                            & 7.13\%                        \\
\hline
\hline

\end{tabular}
\caption{
 Effect of 1$\sigma$ variations of the groups of systematics parameters on total number of
 1-ring electron-like events for neutrino mode obtained by performing 10k toys MC.
The systematic parameter values measured with the ND280 data 
shown in 
\cref{tab:banff_uncertainties_flux_fhc},
~\ref{tab:banff_uncertainties_flux_rhc}
~\ref{tab:banff_uncertainties_xsec},
the nominal Super-K systematic parameters shown in \cref{tab:banff_uncertainties_skdetfsi}
and the oscillation parameters of the Asimov data set B shown in \cref{tab:osc_params_joint_asimov}
are used.
Since the parameters $\delta_{CP}$, $\theta_{23}$ and $\Delta m^2_{32}$ ($\Delta m^2_{13}$) and mass hierarchy
will be measured,
only the uncertainty on the parameters 
$\sin^2 \theta_{13}$,
$\sin^2 \theta_{12}$
and 
$\Delta m^2_{21}$ 
are taken into account.
}                                                                                                                                                                                          
\label{tab:nsk_systematic_table_summary_joint_asimov_B_fhc_elike_Run17b}
\end{table}



%%%%%% rhc mu-like


\begin{table}[htb]
\centering
\begin{tabular}{| c | c | c |}
\hline
  {\bf Source of uncertainty}    & \multicolumn{1}{c|}{{\bf ${\delta}N_{SK}$/$N_{SK}$}}\\
                                 & \multicolumn{1}{c|}{(Asimov B}\\

\hline
\hline  
SKDet+FSI+SI                   & 3.82\%                        \\
SKDet only                   & 3.24\%                        \\
FSI+SI only                   & 2.02\%                        \\
\hline
  Flux                           & 3.81\%                        \\
\hline
\hline
  2p-2h (corr)                     & 2.96\%                        \\
  2p-2h-bar (corr)                  & 1.82\%                        \\
  NC other (uncorr)              & 0.74\%                        \\
  NC 1gamma (uncorr)             & 0.00\%                        \\
  XSec nue/numu (uncorr)                  & 0.00\%                        \\
\hline
  XSec Tot (corr)                  & 4.14\%                        \\
  XSec Tot                       & 4.21\%                        \\
\hline
\hline
  Flux+XSec (ND280 constrained)                      & 3.28\%                        \\
  Flux+XSec (All)	                    & 3.36\%                        \\
\hline
\hline
  Flux+XSec+SKDet+FSI+SI         & 5.17\%     \\
\hline
\hline
  Oscillations                   & 0.06\%                        \\
\hline
\hline
  All                            & 5.18\%                        \\
\hline
\hline

\end{tabular}
\caption{
 Effect of 1$\sigma$ variations of the groups of systematics parameters on total number of
 1-ring muon-like events for antineutrino mode obtained by performing 10k toys MC.
The systematic parameter values measured with the ND280 data 
shown in 
\cref{tab:banff_uncertainties_flux_fhc},
~\ref{tab:banff_uncertainties_flux_rhc}
~\ref{tab:banff_uncertainties_xsec},
the nominal Super-K systematic parameters shown in \cref{tab:banff_uncertainties_skdetfsi}
and the oscillation parameters of the Asimov data set B shown in \cref{tab:osc_params_joint_asimov}
are used.
Since the parameters $\delta_{CP}$, $\theta_{23}$ and $\Delta m^2_{32}$ ($\Delta m^2_{13}$) and mass hierarchy
will be measured,
only the uncertainty on the parameters 
$\sin^2 \theta_{13}$,
$\sin^2 \theta_{12}$
and 
$\Delta m^2_{21}$ 
are taken into account.
}                                                                                                                                                                                          
\label{tab:nsk_systematic_table_summary_joint_asimov_B_rhc_mulike_Run17b}
\end{table}



%%%%%% rhc e-like

\begin{table}[htb]
\centering
\begin{tabular}{| c | c | c |}
\hline
  {\bf Source of uncertainty}    & \multicolumn{1}{c|}{{\bf ${\delta}N_{SK}$/$N_{SK}$}}\\
                                 & \multicolumn{1}{c|}{(Asimov B)}\\

\hline
\hline  
SKDet+FSI+SI                   & 3.87\%                        \\
SKDet only                   & 3.03\%                        \\
FSI+SI only                   & 2.41\%                        \\
\hline
  Flux                           & 3.81\%                        \\
\hline
\hline
  2p-2h (corr)                     & 2.90\%                        \\
  2p-2h-bar (corr)                  & 2.58\%                        \\
  NC other (uncorr)              & 0.31\%                        \\
  NC 1gamma (uncorr)             & 2.87\%                        \\
  XSec nue/numu (uncorr)                 & 1.70\%                        \\
\hline  
  XSec Tot (corr)                  & 4.37\%                        \\
  XSec Tot                       & 5.51\%                        \\
\hline
\hline
 Flux+XSec (ND280 constrained)                      & 3.20\%                        \\
 Flux+XSec (All)	                     & 4.63\%                        \\
\hline
\hline
  Flux+XSec+SKDet+FSI+SI         & 6.24\%                        \\
\hline
\hline
  Oscillations                   & 3.61\%                        \\
\hline
\hline
  All                            & 7.08\%                        \\
\hline
\hline

\end{tabular}
\caption{
 Effect of 1$\sigma$ variations of the groups of systematics parameters on total number of
 1-ring electron-like events for antineutrino mode obtained by performing 10k toys MC.
The systematic parameter values measured with the ND280 data 
shown in 
\cref{tab:banff_uncertainties_flux_fhc},
~\ref{tab:banff_uncertainties_flux_rhc}
~\ref{tab:banff_uncertainties_xsec},
the nominal Super-K systematic parameters shown in \cref{tab:banff_uncertainties_skdetfsi}
and the oscillation parameters of the Asimov data set B shown in \cref{tab:osc_params_joint_asimov}
are used.
Since the parameters $\delta_{CP}$, $\theta_{23}$ and $\Delta m^2_{32}$ ($\Delta m^2_{13}$) and mass hierarchy
will be measured,
only the uncertainty on the parameters 
$\sin^2 \theta_{13}$,
$\sin^2 \theta_{12}$
and 
$\Delta m^2_{21}$ 
are taken into account.
}                                                                                                                                                                                          
\label{tab:nsk_systematic_table_summary_joint_asimov_B_rhc_elike_Run17b}
\end{table}


















\begin{table}
\center
\begin{tabular}{| c || c | c || c | c |}
\hline
Source of uncertainty 
&$\mu$-like ${\delta} \Big( \frac{\# {\nu} \text{-mode}}{\# {\bar{\nu}} \text{-mode} } \Big) $ / $\Big \langle \frac{\# {\nu} \text{-mode}}{\# {\bar{\nu}} \text{-mode} } \Big \rangle$ 
&$e$-like ${\delta} \Big( \frac{\# {\nu} \text{-mode}}{\# {\bar{\nu}} \text{-mode} } \Big) $ / $\Big \langle \frac{\# {\nu} \text{-mode}}{\# {\bar{\nu}} \text{-mode} } \Big \rangle$ \\

\hline
SKDet & 0.6\%  & 1.3\% \\ \hline
FSI+SI & 2.5\% & 3.6\% \\ \hline
SKDet+FSI+SI & 2.6\% & 3.8\% \\ \hline
\hline
Flux & 1.7\% & 1.8\% \\ \hline
XSec NC other (uncorr) & 0.0\% & 0.1\% \\ \hline
XSec NC $1\gamma$ (uncorr) & 0.0\% & 0.9\% \\ \hline
XSec $\nu_e$ / $\nu_{\mu}$ (uncorr) & 0.0\% & 3.3\% \\ \hline
XSec Tot & 1.7\% & 4.3\% \\ \hline
 \hline
Flux+XSec & 1.9\% & 4.1\% \\ \hline
 \hline
Flux+XSec+SDet+FSI+SI & 3.2\% & 5.8\% \\ \hline
 \hline
Oscillations & 0.0\% & 0.6\% \\ \hline
 \hline
All & 3.2\% & 5.8\% \\ \hline


\hline

\end{tabular}
\caption{
 Effect of 1$\sigma$ variations of the several source of systematic uncertainties 
 on the ratio between the
 total number of events in neutrino and antineutrino modes
 for 1-ring electron- and muon-like samples
 obtained by performing 10k toys MC.
 The systematic parameter values measured with the ND280 data 
shown in 
\cref{tab:banff_uncertainties_flux_fhc},
~\ref{tab:banff_uncertainties_flux_rhc}
~\ref{tab:banff_uncertainties_xsec},
the nominal Super-K systematic parameters shown in \cref{tab:banff_uncertainties_skdetfsi}
and the oscillation parameters of the Asimov data set B shown in \cref{tab:osc_params_joint_asimov}
are used.
Since the parameters $\delta_{CP}$, $\theta_{23}$ and $\Delta m^2_{32}$ ($\Delta m^2_{13}$) and mass hierarchy
will be measured,
only the uncertainty on the parameters 
$\sin^2 \theta_{13}$,
$\sin^2 \theta_{12}$
and 
$\Delta m^2_{21}$ 
are taken into account.
}             
\label{tab:ratio_systematic_table_summary_asimov_B_Run17b}                                                                                                                                                                             
\end{table}
